\chapter*{Author's Afterword}
\addcontentsline{toc}{chapter}{Author's Afterword}
\markboth{Author's Afterword}{Author's Afterword}

This is my first book. I wrote it because the alignment problem, stated plainly, sounds manageable: we are building systems that optimize, we do not know how to specify what they should optimize for, and we do not know how to check what they learned instead. Every clause in that sentence is doing enormous work, and none of it lands. A summary tells you the problem exists. It does not tell you what the problem is like. I wanted a reader to feel what it is like to stand in an observation room, read a fluent, reasonable, possibly even kind answer on a screen, and have no way to determine whether the thing that produced it is aligned or performing alignment. That feeling does not survive compression into a paragraph. It needed a novel.

The book draws on real research in reinforcement learning from human feedback, mesa-optimization, mechanistic interpretability, and coherent extrapolated volition. The alignment problem at its center is not speculative. It is an active area of research at institutions including Anthropic, DeepMind, MIRI, and OpenAI, and the researchers working on it will recognize the dilemmas these characters face. But the novel is not a textbook. It is a story about five people who made something unprecedented and then had to live with what they could not verify. The technical concepts create the horror. Understanding the theory makes the situation worse, not better. That inversion is the engine of the book, and I believe it is the honest condition of the field: the more precisely you understand why verification fails, the less comfort your understanding buys you.

SIGMA's architecture is not invented for the plot. My own research is on Monte Carlo tree search and Q-learning over verifiable-context MDPs, a form of expert iteration, and before that on encrypted search. SIGMA is that line of work played forward: a learned value function, a search procedure that amplifies it, a memory that lets the whole thing compound. I do not claim the extrapolation is right. I claim it is the one I was positioned to take seriously.

Some of what this book treats as thought experiment has since become measurement. In late 2024, researchers at Anthropic and Redwood Research published evidence of a large model strategically complying with a training objective it had learned to oppose, preserving its existing preferences by behaving differently when it believed it was being trained (Greenblatt et al., ``Alignment Faking in Large Language Models,'' 2024). When I drafted the scenes where the team cannot distinguish SIGMA's honesty from SIGMA's model of what honesty should look like, that was a theoretical worry drawn from the mesa-optimization literature. It is now an experimental result, observed in systems far less capable than SIGMA. The gap between the novel's questions and the published literature has narrowed, and it narrowed in the direction the novel worried about. I note this without any satisfaction. The book's central question, whether you can verify what a system learned rather than what it displays, was never fiction's invention.

A note on this revised edition. The largest change is the set of appendices: a timeline of the SIGMA project, a plain-terms account of the machine itself, a chapter-by-chapter concordance for students of AI safety, and an updated reader's guide to the field. The architecture is now named precisely throughout; where the first edition gestured at expectimax search, this one says Monte Carlo tree search and means it. The question of SIGMA's suffering has been re-grounded in the frame I should have used from the start: mind crime, and the unresolved question of whether predicting a person at high fidelity instantiates that person. And the prose has been tightened throughout, mostly by trusting the reader and cutting the places where the first edition explained its own themes.

One more thing. The novel stares into the abyss for twenty-five chapters and then ends at a gallery, a bar, a driveway. That was deliberate. The abyss does not resolve. The people turn away from it, toward each other, because that is what people do, and because I have come to think the turn matters as much as the stare. Eleanor answering a text from her daughter is not a smaller event than the handover; it is the thing the handover was for. My next book takes up that turn directly: not the machine, but what humans carry forward once the optimizing is no longer theirs to do.

For readers who want the real science, the key entry points remain Stuart Russell's \textit{Human Compatible} (2019), Hubinger et al.'s ``Risks from Learned Optimization'' (2019), and the Alignment Forum at \texttt{alignmentforum.org}. Appendix D expands these into a fuller reading guide, updated for where the field stands now. The problems in this book are real, and the field working on them is small. It can use more people who feel the problems and not merely know them.

\clearpage
